\chapter{Large-Scale Generative Models}

\section{Diffusion in Practice (e.g., Stable Diffusion)}

Diffusion models have become one of the dominant approaches to large-scale generative modeling, particularly in image synthesis.

\medskip

\noindent
\textbf{From theory to practice.}  
Practical diffusion systems extend the basic framework with architectural and algorithmic improvements:
\begin{itemize}
    \item U-Net architectures for denoising
    \item Conditioning mechanisms (e.g., text embeddings)
    \item Efficient sampling strategies
\end{itemize}

\medskip

\noindent
\textbf{Latent diffusion.}  
Instead of operating directly in pixel space, models operate in a lower-dimensional latent space:
\begin{itemize}
    \item Encode images into latent representations
    \item Perform diffusion in latent space
    \item Decode back to pixel space
\end{itemize}

\medskip

\noindent
\textbf{Advantages.}
\begin{itemize}
    \item Reduced computational cost
    \item Improved scalability
\end{itemize}

\medskip

\noindent
\textbf{Conditioning.}  
Generative models can be guided by additional inputs:
\begin{itemize}
    \item Text prompts
    \item Class labels
    \item Other modalities
\end{itemize}

\medskip

\noindent
\textbf{Insight.}  
Modern diffusion systems combine probabilistic modeling with powerful neural architectures and conditioning mechanisms.

\section{Large Language Models}

Large language models (LLMs) are generative models for sequences, typically based on transformer architectures.

\medskip

\noindent
\textbf{Autoregressive modeling.}  
LLMs model the joint distribution of a sequence:
\[
p(x_1, \dots, x_T) = \prod_{t=1}^T p(x_t \mid x_1, \dots, x_{t-1}).
\]

\medskip

\noindent
\textbf{Training objective.}
\[
\max_\theta \; \mathbb{E} \left[ \sum_{t=1}^T \log p_\theta(x_t \mid x_{<t}) \right].
\]

\medskip

\noindent
\textbf{Transformer backbone.}
\begin{itemize}
    \item Self-attention layers
    \item Feedforward networks
    \item Residual connections and normalization
\end{itemize}

\medskip

\noindent
\textbf{Capabilities.}
\begin{itemize}
    \item Text generation
    \item Question answering
    \item Code generation
    \item Reasoning tasks
\end{itemize}

\medskip

\noindent
\textbf{Pretraining and fine-tuning.}
\begin{itemize}
    \item Pretrained on large corpora
    \item Adapted to downstream tasks
\end{itemize}

\medskip

\noindent
\textbf{Insight.}  
LLMs leverage scale and architecture to model complex sequential dependencies.

\section{Scaling Laws}

Empirical studies reveal systematic relationships between model performance and scale.

\medskip

\noindent
\textbf{Scaling factors.}
\begin{itemize}
    \item Model size (number of parameters)
    \item Dataset size
    \item Compute resources
\end{itemize}

\medskip

\noindent
\textbf{Empirical law.}
\[
\text{Error} \propto N^{-\alpha},
\]
where $N$ represents scale and $\alpha > 0$.

\medskip

\noindent
\textbf{Implications.}
\begin{itemize}
    \item Larger models tend to perform better
    \item Gains diminish with scale
\end{itemize}

\medskip

\noindent
\textbf{Compute-optimal tradeoffs.}
\begin{itemize}
    \item Balance model size and dataset size
    \item Efficient allocation of computational resources
\end{itemize}

\medskip

\noindent
\textbf{Emergent behavior.}
\begin{itemize}
    \item New capabilities appear at large scale
    \item Not easily predicted from small models
\end{itemize}

\medskip

\noindent
\textbf{Insight.}  
Scaling is a fundamental driver of performance in modern machine learning systems.

\section{Engineering Considerations}

Large-scale generative models require careful engineering to train and deploy.

\medskip

\noindent
\textbf{Distributed training.}
\begin{itemize}
    \item Data parallelism
    \item Model parallelism
    \item Pipeline parallelism
\end{itemize}

\medskip

\noindent
\textbf{Hardware acceleration.}
\begin{itemize}
    \item GPUs and specialized accelerators
    \item Memory and bandwidth constraints
\end{itemize}

\medskip

\noindent
\textbf{Optimization challenges.}
\begin{itemize}
    \item Large batch sizes
    \item Learning rate scheduling
    \item Stability at scale
\end{itemize}

\medskip

\noindent
\textbf{Inference considerations.}
\begin{itemize}
    \item Latency vs throughput
    \item Model compression and quantization
\end{itemize}

\medskip

\noindent
\textbf{Data considerations.}
\begin{itemize}
    \item Data quality and diversity
    \item Preprocessing and filtering
\end{itemize}

\medskip

\noindent
\textbf{Safety and alignment.}
\begin{itemize}
    \item Controlling model outputs
    \item Mitigating harmful or biased behavior
\end{itemize}

\medskip

\noindent
\textbf{Insight.}  
Engineering plays a crucial role in transforming theoretical models into practical systems.

\section{A Unifying Perspective}

Large-scale generative models integrate multiple dimensions:

\begin{itemize}
    \item \textbf{Modeling:} probabilistic and neural architectures
    \item \textbf{Computation:} large-scale optimization
    \item \textbf{Data:} massive and diverse datasets
    \item \textbf{Systems:} distributed and efficient implementations
\end{itemize}

\medskip

\noindent
\textbf{Core idea.}  
Performance emerges from the interaction between models, data, and computation.

\medskip

\noindent
\textbf{Key insight.}  
Modern generative models are not defined solely by their mathematical formulation, but by their scale and system-level design.

\section{Outlook}

This chapter examined large-scale generative models:
\begin{itemize}
    \item diffusion models in practice,
    \item large language models,
    \item scaling laws,
    \item engineering considerations.
\end{itemize}

\medskip

This concludes our study of modern generative modeling techniques.

\medskip

\noindent
\textbf{Guiding principle.}  
Advances in machine learning arise from the combination of theory, architecture, data, and computation at scale.